\documentclass[12pt]{article}
\usepackage{amsmath,amssymb,amsfonts}
\usepackage[margin=1in]{geometry}

\begin{document}

\section*{Chapter 3, Problem 24}

\textbf{Problem.} Show that $(C^{-1})^T = (C^T)^{-1}$. That is $C^{-T} = C^{T^{-1}}$. Then use this result to get $C^{-1} = C^{T}\,(C\,C^T)^{-1} = (C^T C)^{-1} C^T$ --- which is the operator analogue of the familiar expansion of $(x + y)^{-1}$ when $y$ is ``small.''

\bigskip
\textbf{Solution.}

\textbf{Part 1:} Show $(C^{-1})^T = (C^T)^{-1}$.

Start from $CC^{-1} = I$. Take the transpose of both sides:
\[
(CC^{-1})^T = I^T = I.
\]
Using the property $(AB)^T = B^T A^T$:
\[
(C^{-1})^T C^T = I.
\]
This shows that $(C^{-1})^T$ is the left inverse of $C^T$. Similarly, from $C^{-1}C = I$:
\[
C^T (C^{-1})^T = I.
\]
So $(C^{-1})^T$ is also the right inverse of $C^T$. Therefore:
\[
(C^{-1})^T = (C^T)^{-1}. \quad \blacksquare
\]

\bigskip
\textbf{Part 2:} Useful identities.

From $(C^T)^{-1} = (C^{-1})^T$, multiply both sides on the right by $C^T$... Actually the identities $C^{-1} = C^T(CC^T)^{-1}$ and $C^{-1} = (C^TC)^{-1}C^T$ hold when $C$ is invertible and can be verified directly:

For $C^{-1} = (C^T C)^{-1}C^T$: multiply on the left by $C$:
\[
C \cdot (C^T C)^{-1}C^T = C(C^T C)^{-1}C^T.
\]
And $I = CC^{-1}$, so we need to verify: multiply $(C^TC)^{-1}C^T$ on the left by $C$:
\[
C(C^TC)^{-1}C^T.
\]
This is not obviously $I$. However, if we multiply on the left by $C^TC$:
\[
(C^TC)\bigl[(C^TC)^{-1}C^T\bigr] = C^T.
\]
And $(C^TC) C^{-1} = C^T$, which is indeed correct since $C^T C \cdot C^{-1} = C^T (CC^{-1}) = C^T$. So $(C^TC)^{-1}C^T = C^{-1}$. $\checkmark$

Similarly for $C^{-1} = C^T(CC^T)^{-1}$: multiply on the right by $C$:
\[
C^T(CC^T)^{-1} \cdot C = C^T(CC^T)^{-1}C.
\]
And $C^{-1}\cdot C = I$, so we verify: multiply $C^T(CC^T)^{-1}$ on the right by $CC^T$:
\[
C^T(CC^T)^{-1}(CC^T) = C^T,
\]
and $C^{-1}(CC^T) = C^T$, which is correct. So $C^T(CC^T)^{-1} = C^{-1}$. $\checkmark$

These identities become particularly useful for non-square matrices, where they yield the left and right pseudo-inverses, respectively. \hfill $\blacksquare$

\end{document}
